\documentclass[12pt, a4paper]{article}

%=============== PAQUETES ===============
\usepackage[utf8]{inputenc}
\usepackage[T1]{fontenc}
\usepackage[spanish,es-noshorthands]{babel}
\usepackage{geometry}
\usepackage{graphicx}
\usepackage{amsmath}
\usepackage{xcolor}
\usepackage{circuitikz}
\usepackage{enumitem}
\usepackage{float}
\usepackage{hyperref}
\usetikzlibrary{arrows.meta, babel}

%=============== CONFIGURACIÓN ===============
\geometry{a4paper, margin=2.5cm}

\hypersetup{
    colorlinks=true,
    linkcolor=blue,
    urlcolor=cyan,
    pdftitle={Circuito de Entrada - SMPS},
}

%=============== DOCUMENTO ===============
\begin{document}

\title{\textbf{Fuente de Alimentación Conmutada (SMPS)}}
\author{Apuntes de Electrónica}
\date{\today}
\maketitle

\section{Circuito de Entrada}
Los componentes de una fuente de alimentación conmutada (SMPS) desde la entrada de la red eléctrica hasta el rectificador primario de onda completa son los siguientes:

\begin{itemize}[itemsep=1.5ex]
    \item \textbf{Fusible:} Es el primer elemento de la entrada y protege al circuito abriéndose si la corriente supera su límite establecido. Debe ser de "fusión lenta" (anti transitorios) porque al conectar la fuente a la red se produce un pico de corriente alto al que debe sobrevivir. A veces, los fabricantes utilizan un \textbf{fusistor}, que es una resistencia especial de muy bajo valor óhmico diseñada para actuar como fusible.
    
    \item \textbf{Filtro EMC (Compatibilidad Electromagnética):} Su propósito es aislar el ruido eléctrico generado por interferencias electromagnéticas, tanto para que no entren desde la red hacia la fuente, como para que la fuente no emita ruido a la instalación eléctrica. Consta principalmente de:
    \begin{itemize}[itemsep=1ex, topsep=1ex]
        \item \textbf{Condensadores de seguridad:} Los de \textbf{tipo X} se conectan entre las dos líneas de alimentación y los de \textbf{tipo Y} se conectan entre la línea y la tierra. Su función es derivar las señales nocivas de alta frecuencia. Son de tipo "autorregenerable", lo que significa que en caso de un arco eléctrico destruyen solo una pequeña zona de su material en lugar de crear un cortocircuito peligroso.
        \item \textbf{Bobinas acopladas (Modo Común):} Normalmente se enrollan sobre un núcleo de ferrita. Se oponen con firmeza al paso de las señales indeseadas de alta frecuencia. Se enrollan deliberadamente en sentidos opuestos para que las tensiones continuas anulen sus campos magnéticos entre sí, evitando que el núcleo de ferrita se sature y pierda sus propiedades magnéticas.
    \end{itemize}

    \item \textbf{Resistencias de descarga:} Son resistencias ubicadas en la entrada (frecuentemente varias en serie) que garantizan la descarga eléctrica de los condensadores cuando se apaga o desconecta el equipo, previniendo posibles descargas al manipular la placa.

    \item \textbf{Interruptor de encendido o puente:} Los terminales de entrada de algunas placas pueden incluir un interruptor (ON/OFF) o bien un puente de alambre soldado para dejar la alimentación conectada permanentemente.

    \item \textbf{Varistor (VDR o MOV):} Es una resistencia dependiente de la tensión, ubicada en paralelo a la entrada. Funciona como protección contra picos transitorios. A tensiones de trabajo normales presenta una resistencia casi infinita (interruptor abierto), pero ante un pico o sobretensión extrema su resistencia cae drásticamente. Esto provoca un cortocircuito intencionado que absorbe el pico y obliga a fundir el fusible, salvando a los transistores y otras partes delicadas del equipo.

    \item \textbf{Termistor NTC:} Es un componente con un coeficiente de temperatura negativo diseñado para frenar el intenso pico de corriente provocado por los condensadores al encender el equipo. Al momento de la conexión la NTC se encuentra fría y ofrece una resistencia elevada limitando el paso de corriente. Inmediatamente después, el propio paso de corriente la calienta, haciendo que su resistencia caiga a un valor casi despreciable y permitiendo el flujo eléctrico normal.

    \item \textbf{Rectificador en puente (Rectificador primario):} Transforma la tensión de corriente alterna de entrada en corriente continua pulsante. Está formado por cuatro diodos que pueden venir en componentes discretos individuales o en un solo encapsulado o pastilla integrada. Al funcionar en conmutación (abriendo y cerrando el paso a la corriente), los diodos generan ruidos y oscilaciones de alta frecuencia, por lo que frecuentemente se coloca una \textbf{red RC (resistencia y condensador en serie)} en paralelo a la salida del rectificador. Esta red ayuda a amortiguar la conmutación y a disipar los armónicos nocivos en forma de calor.
\end{itemize}

\newpage
\subsection{Esquema Eléctrico del Circuito de Entrada}

A continuación, se presenta el diagrama esquemático que ilustra la disposición típica de estos componentes en la etapa de entrada de una SMPS. 

\begin{figure}[H]
    \centering
    \resizebox{\textwidth}{!}{
    \begin{circuitikz}[american, thick, scale=1.0, transform shape]
        
    % Entrada de red (Línea y Neutro)
    \draw (0,4) node[anchor=east] {\textbf{L}} to[short, o-] (0.5,4)
          to[fuse, l=Fusible] (2.5,4) node[circ]{} node[above]{a} -- (3,4);
    \draw (0,0) node[anchor=east] {\textbf{N}} to[short, o-] (3,0);

    % Varistor (VDR) para protección contra sobretensiones
    \draw (2.5,4) to[varistor, l=VDR] (2.5,0);

    % Termistor NTC para limitar la corriente inicial
    \draw (3,4) to[thermistor, l=NTC] (5,4);
    \draw (3,0) to[short] (5,0);

    % Resistencias de descarga
    \draw (5,4) to[short] (6,4);
    \draw (5,0) to[short] (6,0);
    \draw (5.5,4) node[circ]{} node[above]{b} to[R, l=$R_{desc}$] (5.5,0);

    % Filtro EMC - Condensador de seguridad Tipo X
    \draw (6,4) -- (7.5,4) to[C, l=$C_X$] (7.5,0);
    \draw (6,0) -- (7.5,0);

    % Filtro EMC - Bobinas acopladas en modo común
    \draw (7.5,4) -- (8,4) to[L, l=$L_1$] (10,4);
    \draw (7.5,0) -- (8,0) to[L, l=$L_2$] (10,0);
    % Líneas punteadas para indicar el acoplamiento magnético (núcleo)
    \draw[dashed] (8.8,2.1) -- (9.8,2.1);
    \draw[dashed] (8.8,1.9) -- (9.8,1.9);

    % Filtro EMC - Condensadores de seguridad Tipo Y a tierra
    \draw (10,4) to[short] (10.5,4) coordinate(c) node[circ]{} node[above]{c} to[C, l_=$C_{Y1}$] (10.5,2);
    \draw (10,0) to[short] (10.5,0) coordinate(h) node[circ]{} node[below]{d} to[C, l=$C_{Y2}$] (10.5,2);
    \draw (10.5,2) node[circ]{} -- (10.8, 2) node[ground]{}; % Punto de unión y tierra a la derecha

    % Puente rectificador de onda completa (dibujado manualmente con 4 diodos)
    \coordinate (br-left) at (12.0,2);
    \coordinate (br-right) at (15.0,2);
    \coordinate (br-top) at (13.5,3.5);
    \coordinate (br-bottom) at (13.5,0.5);

    % Etiquetas para los nodos del puente
    \node[above left] at (br-left) {h};
    \node[above right] at (br-top) {e};
    \node[right] at (br-right) {f};
    \node[below right] at (br-bottom) {g};

    % Dibujo del puente con los diodos escalados a un tamaño menor
    \begin{scope}
        \ctikzset{diodes/scale=0.6}
        \draw (br-left) to[D, l=$D_1$, *-*] (br-top);
        \draw (br-top) to[D, l=$D_2$, *-*] (br-right);
        \draw (br-left) to[D, l_=$D_3$, *-*] (br-bottom);
        \draw (br-bottom) to[D, l_=$D_4$, *-*] (br-right);
    \end{scope}

    % Conexiones AC desde el filtro al puente rectificador
    \draw (c) -| (br-top);
    \draw (h) -| (br-bottom);

    % Capacitor de filtro entre las salidas de CC (nodos h y f)
    \draw (br-right) |- (16, 5);
    \draw (br-left) |- (16, -1);
    \draw (16, 5) to[pC, l={$68\,\mu\text{F}\text{, } 450\,\text{V}$}, *-*] (16, -1);
    \node[above] at (16, 5) {i};
    \node[below] at (16, -1) {j};

    \end{circuitikz}
    }
    \caption{Esquema integral del circuito de entrada, protección y rectificación de una Fuente Conmutada.}
\end{figure}

\subsection{Análisis Teórico: El Efecto de Tierra Flotante}

Para comprender el comportamiento de las señales antes del puente rectificador a la frecuencia de la red eléctrica (típicamente 50\,Hz o 60\,Hz), supongamos que entre los terminales de entrada \textbf{L} y \textbf{N} se aplica una señal sinusoidal de amplitud $A$, es decir, $V_{in}(t) = A \sin(\omega t)$.

\subsubsection*{1. Comportamiento de los componentes a baja frecuencia}
A la frecuencia de la red eléctrica, el circuito se comporta de una manera muy específica:
\begin{itemize}
    \item \textbf{Bobinas acopladas ($L_1$ y $L_2$):} Su reactancia inductiva ($X_L = 2\pi f L$) a 50/60\,Hz es extremadamente baja. Para efectos prácticos de la señal de red, actúan como un cortocircuito.
    \item \textbf{Termistor NTC y Fusible:} Presentan una resistencia muy baja.
\end{itemize}
En consecuencia, la señal de voltaje que entra entre \textbf{L} y \textbf{N} llega prácticamente intacta a los nodos \textbf{c} y \textbf{d}. Es decir, la diferencia de potencial entre ellos es $V_{cd} \approx A \sin(\omega t)$.

\subsubsection*{2. El divisor de tensión capacitivo ($C_{Y1}$ y $C_{Y2}$)}
El punto clave del análisis radica en el nodo de \textbf{tierra}, el cual está conectado exactamente en el punto medio de los condensadores de seguridad tipo Y ($C_{Y1}$ y $C_{Y2}$). Dado que en un filtro EMC bien diseñado $C_{Y1} = C_{Y2}$, estos actúan como un divisor de tensión capacitivo perfectamente balanceado. 

Lo que se mida respecto a esta ``tierra'' dependerá de cómo esté conectado el equipo a la instalación eléctrica:

\begin{itemize}
    \item \textbf{Escenario A: Equipo sin conexión a tierra (Tierra flotante).} Si medimos respecto al nodo central del divisor (que suele estar conectado al chasis metálico de la fuente), el divisor parte la amplitud de la señal de entrada exactamente por la mitad:
    \begin{itemize}
        \item Señal en el nodo \textbf{c} respecto a tierra: $V_c = \frac{A}{2} \sin(\omega t)$
        \item Señal en el nodo \textbf{d} respecto a tierra: $V_d = -\frac{A}{2} \sin(\omega t)$
    \end{itemize}
    \textit{Nota de ingeniería:} Esta es la razón por la cual, si se toca el chasis metálico de un equipo electrónico con la conexión a tierra interrumpida, se siente un leve hormigueo o descarga. El chasis está flotando a la mitad del voltaje de la red debido a estos condensadores.

    \item \textbf{Escenario B: Equipo con conexión a tierra real.} En la mayoría de las instalaciones eléctricas, el cable Neutro (\textbf{N}) se une físicamente a la toma de Tierra en el cuadro principal. Si este es el caso:
    \begin{itemize}
        \item El nodo \textbf{N} y la tierra están al mismo potencial ($0\,\text{V}$).
        \item Como $L_2$ es casi un cortocircuito a baja frecuencia, la señal en el nodo \textbf{d} respecto a tierra será prácticamente nula: $V_d \approx 0\,\text{V}$.
        \item El nodo \textbf{L} lleva toda la oscilación. Por tanto, la señal en el nodo \textbf{c} respecto a tierra será la onda completa: $V_c \approx A \sin(\omega t)$.
    \end{itemize}
\end{itemize}

\subsubsection*{3. Etapa de Rectificación de Onda Completa}
El puente rectificador está formado por los diodos $D_1$, $D_2$, $D_3$ y $D_4$. Su funcionamiento durante un ciclo completo de la señal alterna que llega desde el filtro EMC es el siguiente:
\begin{itemize}
    \item \textbf{Semiciclo positivo:} Cuando el nodo \textbf{e} es positivo con respecto al nodo \textbf{g}, la corriente fluye a través del diodo $D_2$, pasa por el condensador de filtro en el sentido del nodo \textbf{f} al nodo \textbf{h}, y retorna al nodo \textbf{g} a través del diodo $D_3$.
    \item \textbf{Semiciclo negativo:} Cuando la polaridad se invierte (el nodo \textbf{g} se hace positivo con respecto al nodo \textbf{e}), la corriente fluye a través del diodo $D_4$, atraviesa el condensador exactamente en el mismo sentido que en el caso anterior (de \textbf{f} a \textbf{h}), y retorna al nodo \textbf{e} a través del diodo $D_1$.
\end{itemize}
De esta manera, la corriente circula siempre en una sola dirección, logrando la rectificación de onda completa. El condensador de $68\,\mu\text{F}$ se encarga de suavizar (filtrar) esta tensión continua pulsante, almacenando energía durante los picos y entregándola durante los valles, lo que establece al nodo \textbf{f} como el terminal positivo ($+V_{DC}$) y al nodo \textbf{h} como el retorno negativo (GND o $-V_{DC}$).

\subsubsection*{4. Función de las inductancias $L_1$ y $L_2$ (Filtro de Modo Común)}
Las bobinas $L_1$ y $L_2$ conforman lo que se conoce como un \textbf{choque de modo común} (Common-Mode Choke). Su función principal es suprimir el ruido electromagnético de alta frecuencia generado por la conmutación de la propia fuente, evitando que este contamine la red eléctrica externa, mientras permiten el paso libre de la energía a la frecuencia de red (50/60\,Hz). Su funcionamiento se basa en la dirección de la corriente y el acoplamiento magnético en el núcleo:
\begin{itemize}
    \item \textbf{Señales en Modo Diferencial (Energía de red):} La corriente alterna normal entra por la línea (\textbf{L}), atraviesa $L_1$, alimenta el circuito, y retorna por el neutro (\textbf{N}) atravesando $L_2$ en sentido opuesto. Al estar enrolladas estratégicamente sobre el mismo núcleo de ferrita, los campos magnéticos generados por estas dos corrientes opuestas de baja frecuencia se anulan entre sí. Como resultado, la impedancia magnética neta es casi cero (las bobinas actúan como simples cables) y el núcleo no se satura, sin importar cuánta corriente consuma la fuente.
    \item \textbf{Señales en Modo Común (Ruido electromagnético):} El ruido de alta frecuencia (típicamente entre 10\,kHz y varios MHz), originado por los veloces cambios de encendido y apagado del transistor principal de la SMPS, tiende a propagarse en la misma dirección simultáneamente por ambos cables (\textbf{L} y \textbf{N}) buscando retornar a través de tierra o radiar al ambiente. Para estas corrientes que fluyen en el mismo sentido, los campos magnéticos en $L_1$ y $L_2$ se suman constructivamente. Esto dota a las bobinas de una reactancia inductiva extremadamente alta, actuando como un "muro" o barrera que bloquea de forma efectiva la propagación de estas interferencias.
\end{itemize}

\subsubsection*{5. ¿Forman $R_{desc}$ y $C_X$ un filtro pasa bajo?}
Es una intuición común pensar que la resistencia de descarga ($R_{desc}$) y el condensador $C_X$ conforman un filtro RC pasa bajo. Sin embargo, \textbf{esto es un error conceptual}. 

Para que exista un filtro RC pasa bajo que atenúe el ruido de la línea, la resistencia tendría que estar conectada en \textit{serie} con el flujo de corriente principal. Colocar una resistencia en serie en la entrada de una fuente de alimentación es inviable, ya que toda la corriente de consumo pasaría por ella, disipando una enorme cantidad de energía en forma de calor ($P = I^2R$) y reduciendo drásticamente la eficiencia.

En la realidad, el filtrado EMI de una fuente conmutada no es de tipo RC, sino de tipo \textbf{LC (Inductor-Condensador)}. La atenuación real de alta frecuencia se logra mediante la combinación de las bobinas $L_1$ y $L_2$ (que se oponen a los cambios rápidos de corriente) y los condensadores $C_X$ y $C_{Y}$ (que cortocircuitan el ruido restante).

\textbf{Entonces, ¿cuál es la función de $R_{desc}$?} \\
Dado que $R_{desc}$ está conectada en \textit{paralelo} con $C_X$ y tiene un valor óhmico muy elevado (típicamente del orden de los Megaohmios, $\text{M}\Omega$), su impacto durante el funcionamiento normal es prácticamente nulo y no interviene en el filtrado. Su único y vital propósito es la \textbf{seguridad humana}. 

Cuando el equipo se desconecta del tomacorriente, el condensador $C_X$ queda cargado con la tensión de red. Si no existiera $R_{desc}$, esa carga eléctrica podría permanecer almacenada por horas. Si un usuario tocara las clavijas del enchufe desconectado, recibiría una fuerte descarga. $R_{desc}$ proporciona un camino para que el condensador se descargue de forma segura en una fracción de segundo tras la desconexión de la red.

\subsubsection*{6. Dispositivos de protección: Varistor (VDR) y Termistor NTC}
Además del filtrado, la etapa de entrada cuenta con componentes críticos para proteger la circuitería interna de la fuente frente a anomalías en la red y transitorios mecánicos de encendido:
\begin{itemize}
    \item \textbf{Varistor (VDR o MOV):} Actúa como la primera línea de defensa contra picos de sobretensión transitorios (como los inducidos por descargas atmosféricas o anomalías en la red). Se conecta en paralelo a la entrada. En condiciones normales, su resistencia es casi infinita. Sin embargo, si la tensión supera su umbral de diseño, su resistencia cae abruptamente, absorbiendo (cortocircuitando) la energía del pico transitorio. Si la sobretensión es muy alta o sostenida, este cortocircuito intencional provocará que el fusible principal se funda, desconectando y salvando de forma permanente al resto del equipo.
    \item \textbf{Termistor NTC (Coeficiente de Temperatura Negativo):} Su propósito es proteger a los diodos del puente rectificador frente al enorme pico de corriente inicial de encendido (\textit{inrush current}). Al conectar el equipo a la red, los condensadores principales de filtro (descargados) se comportan por una fracción de segundo como un cortocircuito. El NTC se encuentra frío en el momento de la conexión y presenta una resistencia moderada (por ejemplo, entre $5\,\Omega$ y $10\,\Omega$), lo que frena este pico de corriente. Inmediatamente después, el paso continuo de la corriente de trabajo calienta el termistor, haciendo que su resistencia caiga a un valor cercano a cero ($<1\,\Omega$), permitiendo el flujo normal de energía sin causar caídas de tensión importantes ni penalizar la eficiencia de la fuente.
\end{itemize}

\newpage
\section{Etapa de Conmutación}

Una vez que la tensión de la red eléctrica ha sido rectificada y filtrada en la etapa de entrada, se obtiene una tensión continua (DC) de alto valor (aproximadamente $300\,\text{V}$ a $340\,\text{V}$ en sistemas de $220\,\text{V}_{AC}$, o $150\,\text{V}$ a $170\,\text{V}$ en sistemas de $110\,\text{V}_{AC}$). La \textbf{etapa de conmutación} es el verdadero corazón de la fuente de alimentación conmutada (SMPS). 

Su función principal es "trocear" o conmutar esta tensión continua a una muy alta frecuencia (típicamente entre $20\,\text{kHz}$ y $150\,\text{kHz}$) para aplicarla al devanado primario del transformador. Al trabajar a estas frecuencias, los componentes magnéticos y capacitivos pueden ser drásticamente más pequeños, ligeros y eficientes.

Los componentes principales que conforman esta etapa incluyen:

\begin{itemize}[itemsep=1.5ex]
    \item \textbf{Transistor de Conmutación:} Generalmente se utiliza un MOSFET de potencia (o un BJT en diseños más antiguos). Actúa como un interruptor ultrarrápido que se abre y se cierra miles de veces por segundo, controlando el flujo de corriente hacia el transformador.
    
    \item \textbf{Controlador PWM (Pulse Width Modulation):} Es el "cerebro" lógico de la fuente. Este circuito integrado (como la famosa familia \textbf{UC384x}, por ejemplo, UC3842 o UC3844) genera los pulsos que gobiernan la puerta (\textit{gate}) del MOSFET. Ajustando el ancho de estos pulsos (el tiempo que el transistor permanece encendido frente al tiempo que permanece apagado), el controlador regula la cantidad de energía transferida al secundario y estabiliza la tensión de salida.
    
    \item \textbf{Red Snubber (Amortiguador):} Está compuesta típicamente por un diodo rápido, una resistencia y un condensador de alto voltaje (conocida como red RCD), o en su defecto por un diodo supresor de transitorios (TVS). Se coloca en paralelo con el devanado primario del transformador. Su misión vital es absorber y disipar los letales picos de voltaje inducidos (\textit{flyback spike}) que se generan por la inductancia de dispersión del transformador en el instante exacto en que el transistor corta bruscamente la corriente.
    
    \item \textbf{Transformador de Alta Frecuencia (Chopper):} Transfiere la energía desde el lado primario (alta tensión) al lado secundario (baja tensión) proporcionando \textbf{aislamiento galvánico} por estrictas razones de seguridad. A diferencia de los transformadores convencionales de chapas de acero al silicio diseñados para 50/60\,Hz, el núcleo de este transformador está fabricado de \textbf{ferrita}, un material cerámico magnético ideal para no perder eficiencia a tan altas frecuencias.
\end{itemize}

El diseño coordinado de estos componentes es lo que permite a las SMPS alcanzar eficiencias muy superiores (a menudo $>85\%$) en comparación con las antiguas fuentes de alimentación lineales.

\subsection{Esquema Eléctrico de la Etapa de Conmutación}

A continuación, se presenta el diagrama esquemático de la etapa de conmutación. Para este caso, se ha ilustrado una topología \textbf{Flyback}, la cual es la configuración más utilizada en fuentes de alimentación conmutadas de baja y media potencia (generalmente menores a 150\,W) debido a su bajo costo, simplicidad y bajo número de componentes.

\begin{figure}[H]
    \centering
    \resizebox{0.85\textwidth}{!}{
    \begin{circuitikz}[american, thick, scale=1.0, transform shape, font=\scriptsize]
        
    % Reducimos el tamaño global de los componentes y las etiquetas
    \ctikzset{bipoles/length=0.75cm}
    
    % Estilo para puntos de prueba (Test Points)
    \tikzset{tp/.style={circle, fill=red, inner sep=1.5pt}}
    
    % MOSFET de conmutación
    \draw (4,2) node[nigfete, anchor=D] (mosfet) {};
    \node[below right] at (mosfet.D) {D};
    \node[above right] at (mosfet.S) {S};
    \node[above left] at (mosfet.G) {G};
    
    % Líneas de entrada de voltaje continuo (VDC)
    \draw (-1.5,6) node[ocirc] {} node[left] {i} -- (9,6);
    
    % Diodo entre Source y Drain (ánodo en Source)
    \draw (mosfet.S) node[circ]{} -- ++(1.0,0) node[circ]{} coordinate(tap_S) -- ++(0.5,0) coordinate(aux_S) 
          to[D, l_=$D_1$] (aux_S |- mosfet.D) node[circ]{} -- (mosfet.D);
    
    % Potenciómetro de Puerta (Gate) en la línea j (desplazado a la izquierda)
    \coordinate (pot_start) at (-0.5,0 |- mosfet.G);
    \draw (-1.5,0 |- mosfet.G) node[ocirc] {} node[left] {j} -- (pot_start);
    \draw (pot_start) node[circ] {} to[potentiometer, l=$5.1\,\text{k}\Omega$] (mosfet.G) node[circ] {};
    
    % Tierra analógica entre los nodos i y j (Filtro Y)
    \draw (-1.0, 6) node[circ] {} to[C, l=$C_{Y1}$] (-1.0, 4) node[circ] {} to[C, l=$C_{Y2}$] (-1.0, 0 |- mosfet.G) node[circ] {};
    \draw (-1.0, 4) -- (-1.5, 4) node[ground] {};

    % Diodo Zener en paralelo con el potenciómetro
    \draw (pot_start) -- (-0.5, 2.5) to[zD, l=$D_Z$] (mosfet.G |- 0, 2.5) -- (mosfet.G);
    
    % Arreglo en paralelo (2 Resistencias y 1 Capacitor) desde el nodo i
    \draw (4,6) node[circ] {} to[R, l=$R_{s1}$] (4,4) node[circ] {};
    \draw (5.5,6) node[circ] {} to[R, l_=$R_{s2}$] (5.5,4) node[circ] {};
    \draw (7,6) node[circ] {} to[C, l_=$C_s$] (7,4);
    \draw (4,4) -- (7,4);
    
    % Dos diodos en paralelo con ánodos al Drain del MOSFET
    \draw (4,2) node[circ] {} to[D, l=$D_{s1}$] (4,4);
    \draw (4,2) -- (5.5,2) to[D, l_=$D_{s2}$] (5.5,4);
    
    % Primario del transformador entre el nodo "i" y Drain (D)
    \coordinate (prim_D) at (9, 0 |- mosfet.D);
    \draw (9,6) node[circ] {} to[L, l_=$L_P$] (prim_D) -- (5.5,2);
    \node[circ] at (8.6, 5.6) {}; % Punto de fase (polaridad)
    
    % Núcleo de Ferrita
    \draw[dashed] (9.6, 5.5) -- (9.6, -3);
    \draw[dashed] (9.8, 5.5) -- (9.8, -3);
    
    % Devanado Auxiliar (Otro primario)
    \draw (9, 0) node[circ] {} to[L, l_=$L_{aux}$] (9, -3) node[circ] {};
    \node[circ] at (8.6, -2.6) {}; % Punto de fase abajo
    
    % Rectificación y filtrado del devanado auxiliar
    \draw (9, 0) to[R, l_=$R_{aux}$] (7.8, 0) to[D, l_=$D_{aux}$] (6.6, 0) node[circ] {};
    \draw (6.6, 0) to[pC, l=$C_{aux1}$] (6.6, -3);
    \draw (6.6, 0) -- (5.6, 0) node[circ] (caux2_top) {} to[C, name=Caux2] (5.6, -3);
    \node[below right, xshift=0.2cm] at (Caux2.south east) {$C_{aux2}$};
    \draw (5.6, -3) -- (9, -3);
    \draw (6.6, -3) node[circ] {} node[rground] {};
    
    % Optoacoplador debajo del circuito auxiliar
    \draw[dashed, orange] (6.0, -8.5) rectangle (9.0, -6.5);
    \node[below, font=\scriptsize, text=orange] at (7.5, -8.5) {Optoacoplador};
    
    % Lado primario (Fototransistor)
    \draw (6.8, -7.5) node[npn, xscale=-1] (opto_npn) {};
    \draw (opto_npn.C) |- (6.0, -7.1) node[above left, font=\tiny] {4};
    \draw (opto_npn.E) |- (6.0, -7.9) node[above left, font=\tiny] {3};
    
    % Lado secundario (LED)
    \draw (9.0, -7.1) node[above left, font=\tiny] {1} node[circ] {} -- (8.2, -7.1) to[D] (8.2, -7.9) -- (9.0, -7.9) node[below left, font=\tiny] {2} node[circ] {};
    
    % Resistencias conectadas al ánodo y cátodo del optoacoplador
    \draw (9.0, -7.1) to[R, l=$R_{op1}$] (11.0, -7.1) node[circ] {};
    \draw (9.0, -7.9) to[R, l=$R_{op2}$] (11.0, -7.9) node[circ] {};

    % TL431 (dibujado como SCR) a la derecha del optoacoplador
    \draw (12.5, -8.5) to[thyristor, mirror, l=TL431, n=scr_ctrl] (12.5, -6.5);

    % Arreglo RC (Resistencia y Capacitor en serie) entre compuerta y cátodo del SCR
    \draw (scr_ctrl.G) -- ++(1.5, 0) coordinate (g_node) to[R, l=$R_{gk}$] (g_node |- 0, -5.5) to[C, l=$C_{gk}$] (12.5, -5.5) -- (12.5, -6.5) node[circ] {};

    % Arreglo (Resistencia y Potenciómetro en serie) entre compuerta y ánodo del SCR
    \draw (g_node) node[circ] {} -- ++(1.5, 0) coordinate (g_node2) to[R, l=$R_{ga}$] (g_node2 |- 0, -9.5) to[potentiometer, l=$P_{ga}$] (12.5, -9.5) -- (12.5, -8.5) node[circ] {};

    % Resistencia entre la línea positiva de salida y la compuerta del SCR
    \draw (g_node2) node[circ] {} to[R, l=$R_{sense}$] (g_node2 |- 0, -5.0) node[circ] {};

    % Conexión del cátodo del SCR a la resistencia Rop2
    \draw (12.5, -6.5) -- (11.2, -6.5) |- (11.0, -7.9);

    % Conexión de la resistencia Rop1 al nodo positivo de la etapa de salida
    \draw (11.0, -7.1) -- (11.0, -5.0) -- (21.5, -5.0) |- (20.5, 5.5) node[circ] {};

    % Flechas de luz
    \draw[-stealth] (7.8, -7.3) -- (7.2, -7.4);
    \draw[-stealth] (7.8, -7.7) -- (7.2, -7.6);

    % Controlador PWM (Familia UC384x)
    \draw (-1.0, -3.2) rectangle (2.0, -0.8) node[pos=.5, align=center, yshift=0.5cm] {UC384x};
    % Pines Superiores e Inferiores
    \draw (0.5, -0.8) node[below, font=\tiny] {OUT 6} -- (0.5, -0.5);
    \draw (0.5, -3.2) node[above, font=\tiny] {GND 5} -- (0.5, -3.5) node[rground] {};
    % Pines Izquierdos
    \draw (-1.0, -1.2) node[right, font=\tiny] {8 VREF} -- (-1.3, -1.2);
    \draw (-1.0, -2.0) node[right, font=\tiny] {2 VFB} -- (-1.3, -2.0);
    \draw (-1.0, -2.8) node[right, font=\tiny] {1 COMP} -- (-1.3, -2.8);
    % Pines Derechos
    \draw (2.0, -1.2) node[left, font=\tiny] {VCC 7} -- (2.3, -1.2);
    \draw (2.0, -2.0) node[left, font=\tiny] {Isense 3} -- (2.3, -2.0);
    \draw (2.0, -2.8) node[left, font=\tiny] {RT/CT 4} -- (2.3, -2.8);
    
    % Capacitor de temporización desde el pin 4 (RT/CT) a tierra
    \draw (2.3, -2.8) -- (2.3, -4.0) -- (1.2, -4.0) -- (1.2, -8.5) to[C, l=$C_t$] (1.2, -9.5) node[rground] {};

    % Resistencia de temporización entre VREF (Pin 8) y RT/CT (Pin 4)
    \draw (-4.5, -5.5) node[circ] {} to[R, l=$R_t$] (1.2, -5.5) node[circ] {};

    % Red de compensación entre VFB y COMP (Resistencia y Capacitor en paralelo)
    \draw (-1.3, -2.0) -- (-2.2, -2.0) node[circ] {};
    \draw (-1.3, -2.8) -- (-2.2, -2.8) node[circ] {};
    \draw (-2.2, -2.0) to[C, l=$C_{comp}$] (-2.2, -2.8);
    \draw (-2.2, -2.0) -- (-3.0, -2.0) node[circ] {} to[R, l_=$R_{comp}$] (-3.0, -2.8) node[circ] {} -- (-2.2, -2.8);

    % Resistencias en serie desde el nodo inferior de Rcomp hacia tierra digital
    \draw (-3.0, -2.8) to[R, l_=$R_{fb1}$] (-3.0, -3.8) node[circ] {} to[R, l_=$R_{fb2}$] (-3.0, -4.8) node[rground] {};

    % Conexión desde el nodo entre Rfb1 y Rfb2 al pin 3 del optoacoplador
    \draw (-3.0, -3.8) -- (-4.0, -3.8) |- (6.0, -7.9);

    % Conexión desde el pin 8 (VREF) al pin 4 del optoacoplador
    \draw (-1.3, -1.2) -- (-4.5, -1.2) node[circ] (vref_node) {} |- (6.0, -7.1);

    % Capacitor a tierra digital desde el nodo VREF
    \draw (vref_node) -- (-5.5, -1.2) node[circ] (vref_c_node) {} -- (-5.5, -8.5) to[C, l_=$C_{vref}$] (-5.5, -9.5) node[rground] {};

    % Transistor PNP encima del capacitor Cvref (Emisor abajo, Colector arriba)
    \draw (-5.5, 0.5) node[pnp, yscale=-1] (pnp_ctrl) {};
    \node[right] at (-5.3, 0.5) {$T_1$};

    % Capacitor entre el colector y la base del transistor PNP
    \draw (pnp_ctrl.C) |- (-7.5, 1.5) node[circ] (c_pnp_node) {} to[C, l_=$C_{cb}$] (-7.5, 0.5) node[circ] {} -- (pnp_ctrl.B);

    % Tierra digital a partir del nodo superior izquierdo
    \draw (c_pnp_node) -- ++(-1.0, 0) node[rground] {};

    % Resistencia desde la base del PNP hacia VREF
    \draw (pnp_ctrl.B) node[circ] {} -- (-7.5, 0.5) to[R, l=$R_b$] (-7.5, -1.2) node[circ] (vref_rb_node) {} -- (vref_c_node);

    % Transistor NPN en la zona por debajo del PNP (Emisor y Colector a la izquierda)
    \draw (-7.5, -4.0) node[npn, xscale=-1] (npn_ctrl) {};
    \node[left] at (-7.7, -4.0) {$T_2$};

    % Conexión de la base del NPN al pin 4 (RT/CT)
    \draw (npn_ctrl.B) -- (-6.0, -4.0) |- (1.2, -6.0) node[circ] {};

    % Resistencia desde el emisor del transistor NPN hacia el pin 3 (Isense)
    \draw (npn_ctrl.E) |- (-2.0, -6.6) to[R, l=$R_e$] (2.6, -6.6) node[circ] {};

    % Conexión del colector del NPN hacia VREF
    \draw (npn_ctrl.C) -- (vref_rb_node);

    % Conexión del emisor del transistor PNP al pin 1 (COMP)
    \draw (pnp_ctrl.E) -- (-5.5, -0.4) -| (-4.2, -2.8) -- (-3.0, -2.8);

    % Línea desde el Gate del MOSFET hasta el pin OUT del PWM con red en paralelo
    \coordinate (g_top) at (mosfet.G |- 0, 0.6);
    \coordinate (g_bot) at (mosfet.G |- 0, -0.2);
    \draw (mosfet.G) -- (g_top) node[circ] {};
    \draw (g_top) -- ++(-0.4,0) coordinate (gt_L) to[R, l_=$R_g$] (gt_L |- g_bot) -- (g_bot);
    \draw (g_top) -- ++(0.4,0) coordinate (gt_R) to[D, l=$D_g$] (gt_R |- g_bot) -- (g_bot);
    \draw (g_bot) node[circ] {} |- (0.5, -0.5);
    
    % Conexión del pin VCC al nodo i
    \draw (2.3, -1.2) -- (2.3, -0.2) to[R, l=$R_{vcc}$] (2.3, 0.8) node[circ] {} -- (2.3, 3) to[R, l=$R_{start1}$] (2.3, 4.5) to[R, l=$R_{start2}$] (2.3, 6) node[circ] {};
    \draw (1.0, -0.2) node[rground] {} to[zD, l=$D_{Z2}$] (1.0, 0.8) -- (2.3, 0.8);
    
    % Conexión del pin 7 (VCC) al nodo sobre Caux2
    \draw (caux2_top) -- (5.2, 0) |- (2.3, -1.2) node[circ] {};

    % Conexión del pin Isense al Source del MOSFET
    \coordinate (isense_node) at (4, -2.0);
    \draw (isense_node) node[circ]{} to[R, l=$R_{isense}$] (2.6, -2.0) node[circ] (isense_filter) {} -- (2.0, -2.0);
    \draw (tap_S) |- (isense_node);
    
    % Capacitor de filtro Isense a tierra digital
    \draw (isense_filter) -- (2.6, -8.5) to[C, l_=$C_f$] (2.6, -9.5) node[rground]{};

    % Dos resistencias en paralelo desde la línea Isense hacia tierra digital
    \draw (isense_node) -- (4.0, -3.8) node[circ]{};
    \draw (3.6, -3.8) -- (4.4, -3.8);
    \draw (3.6, -3.8) to[R, l_=\rotatebox{90}{$R_{cs1}$}] (3.6, -5.3);
    \draw (4.4, -3.8) to[R, l_=\rotatebox{90}{$R_{cs2}$}] (4.4, -5.3);
    \draw (3.6, -5.3) -- (4.4, -5.3);
    \draw (4.0, -5.3) node[circ]{} node[rground]{};

    % Secundario del transformador
    \draw (10.5, 5.5) node[circ] (ls_top) {} to[L, l=$L_S$] (10.5, 2.5) node[circ] (ls_bot) {};
    \node[circ] at (10.9, 2.9) {}; % Punto de fase (inversor)
    
    % Diodos Schottky en paralelo rectificando la salida
    \draw (ls_top) -- (11.2, 5.5) node[circ] {} -- (11.2, 6.2) to[sD, l=$D_{out1}$] (12.8, 6.2) -- (12.8, 5.5) node[circ] {};
    \draw (11.2, 5.5) -- (11.2, 4.8) node[circ] {} to[sD, l_=$D_{out2}$] (12.8, 4.8) node[circ] {} -- (12.8, 5.5);
    \draw (12.8, 5.5) -- (17.5, 5.5) to[L, l=$L_{out}$] (19.5, 5.5) node[circ] {} -- (20.5, 5.5) coordinate (v_out_pos);

    % Línea de retorno del secundario (negativo)
    \draw (ls_bot) -- (20.5, 2.5) node[rground] {};

    % Resistencia de descarga a la salida
    \draw (19.5, 5.5) to[R, l=$R_{out}$] (19.5, 2.5) node[circ] {};

    % Indicador de salida (Resistencia en serie con diodo LED)
    \draw (20.5, 5.5) to[R, l=$R_{led}$] (20.5, 4.0) to[led, l=LED] (20.5, 2.5);

    % Terminales de salida principal (+ y -)
    \draw (21.5, 5.5) node[circ] {} -- (22.5, 5.5) node[ocirc] {} node[right, font=\large\bfseries] {$+$};
    \draw (20.5, 2.5) node[circ] {} -- (22.5, 2.5) node[ocirc] {} node[right, font=\large\bfseries] {$-$};

    % Red Snubber (RC) en paralelo con los diodos
    \draw (11.2, 4.8) -- (11.2, 3.8) to[R, l_=$R_{sn}$] (12.0, 3.8) to[C, l_=$C_{sn}$] (12.8, 3.8) -- (12.8, 4.8);

    % 3 Capacitores electrolíticos de filtro en paralelo a la salida
    \draw (14.0, 5.5) node[circ] {} to[pC, l=$C_{out1}$] (14.0, 2.5) node[circ] {};
    \draw (15.5, 5.5) node[circ] {} to[pC, l=$C_{out2}$] (15.5, 2.5) node[circ] {};
    \draw (17.0, 5.5) node[circ] {} to[pC, l=$C_{out3}$] (17.0, 2.5) node[circ] {};

    % Capacitor desde el extremo inferior de Ls a tierra analógica
    \draw (10.5, 2.5) to[C, l=$C_{ls}$] (10.5, 0.5) node[ground] {};

    % Conexión de la tierra digital del circuito auxiliar al nodo negativo de la etapa de salida
    \draw (9, -3) -- (9, -4.5) to[C, l=$C_Y$] (10.0, -4.5) -| (11.5, 2.5) node[circ]{};

    % Recuadro punteado para la etapa de salida
    \draw[dashed, blue] (10.1, -0.2) rectangle (23.5, 7.0);
    \node[above, font=\scriptsize, text=blue] at (16.8, 7.0) {Etapa de Salida};

    % Recuadro punteado para la etapa de control
    \draw[dashed, red] (-9.2, -10.4) -- (4.8, -10.4) -- (4.8, -1.0) -- (3.8, -1.0) -- (3.8, -0.6) -- (-3.0, -0.6) -- (-3.0, 2.2) -- (-9.2, 2.2) -- cycle;
    \node[above, font=\scriptsize, text=red] at (-6.1, 2.2) {Etapa de Control};

    % --- Puntos de Prueba (Troubleshooting) ---
    % Lado Primario
    \node[tp, label={[label distance=-2pt]90:TP1}] at (9,6) {}; % Vbus+
    \node[tp, label={[label distance=-2pt]270:TP2}] at (mosfet.S) {}; % Vbus- (Source del MOSFET)
    \node[tp, label={[label distance=-2pt]225:TP3}] at (mosfet.G) {}; % Gate Drive
    \node[tp, label={[label distance=-2pt]45:TP4}] at (mosfet.D) {}; % Drain (Switching Node)
    \node[tp, label={[label distance=-2pt]90:TP5}] at (isense_node) {}; % Isense
    \node[tp, label={[label distance=-2pt]270:TP6}] at (2.3, -1.2) {}; % VCC
    \node[tp, label={[label distance=-2pt]135:TP7}] at (vref_node) {}; % VREF
    % Lado Secundario
    \node[tp, label={[label distance=-2pt]90:TP8}] at (scr_ctrl.G) {}; % Feedback (TL431 Ref)
    \node[tp, label={[label distance=-2pt]45:TP9}] at (12.8, 5.5) {}; % Salida rectificada (pre-filtro)
    \node[tp, label={[label distance=-2pt]270:TP10}] at (22.5, 5.5) {}; % Salida final +
    \node[tp, label={[label distance=-2pt]90:TP11}] at (22.5, 2.5) {}; % Salida final -

    \end{circuitikz}
    }
    \caption{Esquema eléctrico básico de la etapa de conmutación (Topología Flyback).}
\end{figure}

\subsection{Análisis Teórico: El Rol del MOSFET y el Transformador Chopper}

Para comprender a fondo la etapa de conmutación en una topología Flyback, es vital analizar la interacción dinámica entre sus dos protagonistas: el transistor MOSFET y el transformador de alta frecuencia (Chopper). A diferencia de las fuentes lineales, donde los componentes operan en un régimen continuo, aquí el funcionamiento se basa en la acumulación y liberación de energía en ciclos de encendido y apagado.

\subsubsection*{1. El MOSFET como Interruptor de Alta Velocidad}
El MOSFET de potencia no opera en su región lineal (como un amplificador), sino exclusivamente en corte y saturación, actuando como un interruptor ideal controlado por voltaje:
\begin{itemize}
    \item \textbf{Estado de Conducción (ON):} Cuando el controlador PWM aplica un pulso positivo en la compuerta (Gate), el MOSFET se satura y cierra el circuito. Esto permite que una corriente creciente circule desde la fuente de alto voltaje continuo (nodo "i") a través del devanado primario del transformador ($L_P$) hacia tierra.
    \item \textbf{Estado de Bloqueo (OFF):} Al retirar el voltaje de la compuerta, el MOSFET se abre abruptamente, interrumpiendo el flujo de corriente. La rapidez con la que el MOSFET puede conmutar entre estos dos estados minimiza las pérdidas por conmutación térmica y permite operar a decenas o cientos de kilohercios.
\end{itemize}

\subsubsection*{2. El Transformador Chopper: Un Inductor Acoplado}
En la topología Flyback, el transformador no actúa como un transformador clásico que transfiere energía simultáneamente del primario al secundario. Su comportamiento real es el de un \textbf{inductor doblemente acoplado}:
\begin{itemize}
    \item \textbf{Fase de Almacenamiento (MOSFET ON):} Mientras el MOSFET conduce, la corriente fluye por el primario ($L_P$). Por convención de los puntos de fase (puntos en los devanados), la polaridad inducida en el secundario ($L_S$) hace que el diodo rectificador de salida quede polarizado en inversa. Por tanto, no hay transferencia de energía al secundario; toda la energía eléctrica de la red se almacena en forma de campo magnético en el entrehierro (\textit{gap}) del núcleo de ferrita.
    \item \textbf{Fase de Transferencia (MOSFET OFF):} Cuando el MOSFET se apaga, el campo magnético acumulado colapsa repentinamente. Por la Ley de Lenz, los devanados invierten su polaridad para intentar mantener la corriente. Esta inversión polariza en directa el diodo del secundario, permitiendo que toda la energía magnética almacenada se vacíe abruptamente hacia los condensadores de filtro y la carga en forma de un pulso de corriente.
\end{itemize}

Además de esta función de almacenamiento y transferencia por "bombeo", el transformador Chopper cumple otras dos misiones críticas: proporciona \textbf{aislamiento galvánico} total entre la peligrosa tensión de la red eléctrica y la salida segura de bajo voltaje, y adapta el nivel de tensión final según su relación de espiras.

\subsubsection*{3. El Significado de los Puntos de Fase Invertidos}
En el esquema de la etapa de conmutación, se observan unos pequeños círculos negros o "puntos de fase" junto a los devanados del transformador. Estos puntos indican la polaridad relativa instantánea de las bobinas. Por convención, si una corriente entra por el terminal con punto en el primario, el terminal con punto en el secundario tendrá un potencial positivo respecto a su otro extremo.

En la topología Flyback, es fundamental notar que el punto del secundario ($L_S$) está \textbf{invertido} eléctricamente respecto al primario ($L_P$) y en contra de la dirección de conducción del diodo rectificador de salida. Esto se diseña así con un propósito específico:
\begin{itemize}
    \item \textbf{Bloqueo intencional:} Cuando el MOSFET se enciende e inyecta un pulso positivo en el punto del primario, el punto del secundario también se vuelve positivo. Sin embargo, como el diodo Schottky está conectado en la terminal opuesta (sin punto), su ánodo recibe un potencial negativo frente a su cátodo. El diodo queda \textbf{polarizado en inversa} e impide que la energía pase a la carga, forzando al transformador a guardar la energía en su campo magnético.
    \item \textbf{Liberación de energía:} Al apagarse el MOSFET, el campo magnético colapsa y las polaridades se invierten por la Ley de Lenz. El terminal \textit{sin punto} del secundario se vuelve fuertemente positivo, polarizando ahora sí en directa al diodo y permitiendo que toda la energía magnética almacenada salga hacia los capacitores de filtro como un pulso de corriente continua.
\end{itemize}
Si los devanados no estuvieran invertidos, el circuito intentaría transferir energía directamente mientras el MOSFET conduce. Eso corresponde a una topología distinta conocida como \textit{Forward}, la cual requiere componentes magnéticos adicionales a la salida.

\subsubsection*{4. La Red Snubber (RCD) como Protección Vital}
Al observar el primario del transformador en el esquema, se aprecia una red compuesta por diodos rápidos ($D_{s1}, D_{s2}$), resistencias ($R_{s1}, R_{s2}$) y un capacitor ($C_s$) conectados en paralelo al devanado $L_P$. Esta configuración se conoce como \textbf{Red Snubber RCD} (Resistencia-Capacitor-Diodo) y es absolutamente crítica para la supervivencia del circuito:
\begin{itemize}
    \item \textbf{El problema de la inductancia de dispersión:} En el mundo real, los transformadores no son perfectos. Una pequeña porción del campo magnético del primario no logra acoplarse al secundario y se queda "atrapada" localmente en el devanado. Esto se modela como un pequeño inductor en serie llamado \textit{inductancia de dispersión} (Leakage Inductance). Cuando el MOSFET se apaga bruscamente, cortando la corriente en nanosegundos, esta inductancia se opone al cambio repentino, generando un gigantesco pico de voltaje transitorio ($V = L \frac{di}{dt}$).
    \item \textbf{El riesgo:} Sin protección, este agresivo pico de voltaje de alta frecuencia se sumaría a la tensión continua de entrada ($V_{DC}$) y a la tensión reflejada del secundario, superando rápidamente el voltaje máximo de ruptura estipulado entre Drenador y Surtidor ($V_{DSS}$) del MOSFET, lo que lo destruiría por avalancha en cuestión de milisegundos.
    \item \textbf{El funcionamiento del Snubber:} 
    \begin{itemize}[itemsep=1ex, topsep=1ex]
        \item \textbf{Absorción (MOSFET se apaga):} En el instante en que el MOSFET corta y el voltaje en su Drenador se dispara de forma letal, los diodos de la red Snubber ($D_{s1}, D_{s2}$) se polarizan en directa. Esto proporciona un camino alternativo seguro para la corriente de la inductancia de dispersión, la cual fluye hacia el capacitor $C_s$, "recortando" (\textit{clamping}) el pico de voltaje a un nivel manejable para el transistor.
        \item \textbf{Disipación (MOSFET conduce):} Durante el siguiente ciclo, cuando el MOSFET se enciende nuevamente, los diodos del Snubber quedan polarizados en inversa (bloqueando el riel principal). En ese momento, la energía destructiva que el capacitor $C_s$ atrapó se descarga lenta y controladamente a través de las resistencias de purga ($R_{s1}, R_{s2}$), disipándose de forma segura como calor y dejando al capacitor listo para absorber el impacto del siguiente ciclo de conmutación.
    \end{itemize}
\end{itemize}

\subsubsection*{5. Arranque (Kick-start) y Auto-alimentación (Devanado Auxiliar)}
El funcionamiento continuo del controlador PWM (familia UC384x) requiere una fuente de alimentación estable en su pin $V_{CC}$ (Pin 7). Sin embargo, reducir directamente los $\approx 310\,\text{V}_{DC}$ del bus principal a los $12\,\text{V}-15\,\text{V}$ necesarios para el chip mediante resistencias es extremadamente ineficiente. Por ello, se utiliza un ingenioso mecanismo de dos fases:
\begin{itemize}
    \item \textbf{Fase de Arranque (Kick-start):} Al conectar la fuente a la red eléctrica, el controlador PWM está apagado. Una pequeña "corriente de goteo" fluye desde el bus de alto voltaje (nodo "i") a través de las resistencias de arranque de alto valor óhmico ($R_{start1}$ y $R_{start2}$). Esta corriente es diminuta pero suficiente para cargar lentamente los capacitores asociados al pin $V_{CC}$ (principalmente $C_{aux2}$). Cuando la tensión acumulada alcanza el umbral de encendido del circuito integrado (típicamente unos $16\,\text{V}$ para el UC3842), el controlador "despierta" y envía el primer pulso al MOSFET.
    \item \textbf{El Problema del Consumo:} Una vez que el PWM comienza a conmutar el MOSFET a alta velocidad, su consumo de corriente se dispara. Las resistencias de arranque no pueden suministrar esta nueva demanda; si se redujera su valor óhmico para dejar pasar más corriente, disiparían un calor destructivo e inaceptable ($P = V^2/R$).
    \item \textbf{Fase de Auto-alimentación:} Apenas se inicia la primera conmutación del MOSFET, el transformador Chopper entra en operación. El devanado auxiliar ($L_{aux}$), que actúa como un pequeño secundario alojado en el propio lado primario, capta la fluctuación magnética y genera una tensión alterna de alta frecuencia. Esta tensión fluye a través de $R_{aux}$ (limitadora de picos), es rectificada por el diodo rápido $D_{aux}$ y filtrada firmemente por $C_{aux1}$ y $C_{aux2}$, convirtiéndose en un suministro DC robusto. A partir de este momento, el circuito provee su propia energía (\textit{Bootstrap}) y la escasa aportación que sigue entregando la rama de arranque inicial se vuelve insignificante, logrando así una alta eficiencia.
\end{itemize}

\newpage
\section{Etapa de Salida y Realimentación}

La etapa de salida y realimentación es la responsable de convertir la energía transferida por el transformador de alta frecuencia (Chopper) nuevamente en corriente continua (CC) estable y limpia de ruidos, así como de informar al controlador primario sobre el estado del voltaje para mantener la regulación.

\subsection{Rectificación y Filtro de Salida (Filtro LC)}
Una vez que la energía cruza el aislamiento galvánico a través del transformador, el devanado secundario entrega una tensión alterna de alta frecuencia (típicamente de forma de onda cuadrada o trapezoidal). Para convertirla en CC útil, se emplean los siguientes componentes:

\begin{itemize}[itemsep=1.5ex]
    \item \textbf{Diodos Rectificadores (Schottky):} A diferencia de los rectificadores de red a baja frecuencia (como el 1N4007), en la salida del secundario se emplean diodos de barrera Schottky o de recuperación ultrarrápida. Esto es vital debido a la alta frecuencia de conmutación. Frecuentemente se colocan varios en paralelo ($D_{out1}$ y $D_{out2}$) para compartir la carga de corriente, reducir el estrés térmico y minimizar la caída de tensión directa, lo que aumenta la eficiencia global.
    \item \textbf{Red Snubber Secundaria:} Formada por una resistencia ($R_{sn}$) y un condensador ($C_{sn}$) en serie, colocados en paralelo con los diodos rectificadores. Su función es absorber los picos de voltaje transitorios (\textit{ringing}) generados por la inductancia parásita del transformador y la capacitancia de los diodos al pasar al estado de bloqueo.
    \item \textbf{Filtro LC (Inductor y Capacitores):} La tensión rectificada por los diodos aún contiene un fuerte rizado (\textit{ripple}) de alta frecuencia que debe eliminarse.
    \begin{itemize}[itemsep=1ex, topsep=1ex]
        \item \textbf{Banco de Capacitores:} Se utilizan varios condensadores electrolíticos en paralelo ($C_{out1}, C_{out2}, C_{out3}$). Esto reduce drásticamente la Resistencia Serie Equivalente (ESR) total del arreglo, mejorando el filtrado y prolongando la vida útil de los componentes al evitar el sobrecalentamiento.
        \item \textbf{Inductor de Salida ($L_{out}$):} Colocado en serie con la salida, forma un robusto filtro pasabajos de tipo LC (o filtro Pi si se añade otro capacitor al final). La bobina se opone a los cambios bruscos de corriente, integrando los pulsos y suavizando profundamente la entrega de energía hacia la carga.
    \end{itemize}
    \item \textbf{Resistencia de Carga Mínima (Bleeder) e Indicador:} La resistencia $R_{out}$ asegura que los condensadores se descarguen progresivamente cuando la fuente se apaga, previniendo descargas accidentales. Adicionalmente, la rama con $R_{led}$ proporciona un indicador visual ("Power Good").
\end{itemize}

\subsection{El Lazo de Realimentación (Feedback) con TL431}
Para que la fuente mantenga un voltaje de salida constante sin importar las variaciones en la carga conectada o en el voltaje de la red eléctrica, se implementa un sistema de control de lazo cerrado (\textit{Closed Loop}). 

\begin{itemize}[itemsep=1.5ex]
    \item \textbf{El Regulador TL431:} Es el corazón del circuito de muestreo en el secundario. Aunque a menudo se le representa y conoce como un diodo Zener programable, internamente funciona como un \textbf{amplificador de error} de alta precisión. Cuenta con un pin de referencia (REF) que posee un umbral de activación interno exacto de $2.5\,\text{V}$. 
    
    \item \textbf{Divisor de Tensión (Sensado):} Se conecta una red de resistencias actuando como divisor de tensión desde la línea positiva de la salida regulada hacia el pin REF del TL431. Los valores de estas resistencias se calculan para que, exactamente cuando la fuente alcanza su voltaje nominal de diseño (por ejemplo, $12\,\text{V}$ o $24\,\text{V}$), el voltaje de la muestra en el pin REF sea exactamente $2.5\,\text{V}$.
    
    \item \textbf{Interacción con el Optoacoplador:} El cátodo del TL431 se conecta en serie con el LED emisor de luz interno del \textbf{optoacoplador}. 
    \begin{itemize}[itemsep=1ex, topsep=1ex]
        \item Si el voltaje de salida \textbf{sube} por encima del valor deseado (por ejemplo, al desconectar una carga pesada), el divisor de tensión eleva el pin REF del TL431 por encima de los $2.5\,\text{V}$. El TL431 responde conduciendo más corriente del cátodo al ánodo, lo que hace que el LED del optoacoplador brille con mayor intensidad.
        \item El fototransistor ubicado en el lado primario capta esta luz intensa, conduce más corriente, y altera el voltaje en el pin de Compensación (COMP) del controlador PWM (ej. UC384x).
        \item El controlador PWM interpreta esta señal y, de manera instantánea, \textbf{reduce el ciclo de trabajo (\textit{Duty Cycle})} de los pulsos que envía a la compuerta (Gate) del MOSFET principal. Al estar menos tiempo encendido, el transformador transfiere menos energía al secundario, provocando que el voltaje de salida baje hasta estabilizarse nuevamente.
        \item Si el voltaje de salida \textbf{baja} (por aumento de consumo de la carga), ocurre el proceso inverso: el pin REF recibe menos de $2.5\,\text{V}$, el TL431 restringe el paso de corriente, el optoacoplador se oscurece, y el controlador PWM compensa \textbf{aumentando el ciclo de trabajo}.
    \end{itemize}
\end{itemize}

Gracias a esta dinámica milimétrica entre el amplificador de error TL431 y el optoacoplador, la fuente de alimentación conmutada logra una regulación extremadamente precisa y continua. Además, garantiza el \textbf{aislamiento galvánico} de seguridad estricto entre la zona de conmutación de alta tensión (primario) y la zona de bajo voltaje que manipulará el usuario (secundario).

\newpage
\section{Ejemplo de Diseño: Fuente de 12\,V y 10\,A (120\,W)}

A continuación se presenta una tabla con valores de componentes sugeridos para dimensionar el circuito previamente analizado, con el objetivo de suministrar una tensión de salida regulada de $\mathbf{12\,\text{V}}$ y una corriente de hasta $\mathbf{10\,\text{A}}$ ($120\,\text{W}$ de potencia). Estos valores asumen una tensión de entrada universal de red eléctrica ($85 - 265\,\text{V}_{AC}$) operando bajo la topología Flyback.

\begin{table}[H]
    \centering
    \caption{Lista de materiales (BOM) sugerida para una SMPS de $12\,\text{V} / 10\,\text{A}$.}
    \renewcommand{\arraystretch}{1.3}
    \begin{tabular}{|p{3.5cm}|p{4.0cm}|p{7.5cm}|}
        \hline
        \textbf{Referencia} & \textbf{Componente} & \textbf{Especificaciones y Notas del Diseño} \\ \hline
        MOSFET ($T_1$) & \textbf{STP11NK90Z} & Transistor Canal N, $900\,\text{V}$, $10\,\text{A}$. Su alta tensión de ruptura soporta el pico de voltaje reflejado del transformador. \\ \hline
        Controlador PWM & \textbf{UC3842} o \textbf{UC3843} & Controlador en modo corriente clásico y robusto. \\ \hline
        Transformador & \textbf{Núcleo ETD34} o EE35 & Frecuencia $\approx 65\,\text{kHz}$. Requiere entrehierro (gap). El secundario debe bobinarse con alambre Litz o hilos múltiples para soportar $10\,\text{A}$. \\ \hline
        Sensado ($R_{isense}$) & \textbf{$0.33\,\Omega$ / $2\,\text{W}$} & Limita la corriente pico primaria a aprox. $3\,\text{A}$ ($V_{sense(max)}=1\,\text{V}$). \\ \hline
        Diodos Snubber \newline ($D_{s1}, D_{s2}$) & \textbf{UF4007} o \textbf{FR107} & Diodos de recuperación ultrarrápida, $1000\,\text{V}$. Jamás usar 1N4007 aquí por ser muy lentos. \\ \hline
        Red Snubber \newline ($R_{s1,2}, C_s$) & \textbf{$47\,\text{k}\Omega$ / $2\,\text{W}$} \newline \textbf{$2.2\,\text{nF}$ / $1\,\text{kV}$} & RC configurado para absorber los transitorios (\textit{spikes}) mortales de la inductancia de dispersión. \\ \hline
        Rectificador Sec. \newline ($D_{out1}, D_{out2}$) & \textbf{MBR2060CT} & Diodo Schottky doble en encapsulado TO-220, $60\,\text{V}$, $20\,\text{A}$. Su baja caída de tensión directa mejora la eficiencia térmica. \\ \hline
        Filtro de Salida \newline ($C_{out1,2,3}$) & $\mathbf{3 \times 1000\,\mu\text{F}}$ & Capacitores electrolíticos especiales de baja ESR (\textit{Low-ESR}), de $25\,\text{V}$. Al estar en paralelo soportan mejor el fuerte rizado de $10\,\text{A}$. \\ \hline
        Filtro LC ($L_{out}$) & \textbf{$4.7\,\mu\text{H}$ / $10\,\text{A}$} & Inductor de choque o varilla con núcleo de polvo de hierro. \\ \hline
        Regulador y Opto. & \textbf{TL431} y \textbf{PC817} & Componentes estándar para el lazo de realimentación. \\ \hline
        Divisor de Voltaje \newline ($R_{sense}, R_{ga}, P_{ga}$) & $R_{sense} = \mathbf{39\,\text{k}\Omega}$ \newline $R_{ga} = \mathbf{9.1\,\text{k}\Omega}$ \newline $P_{ga} = \mathbf{2\,\text{k}\Omega}$ (Pot.) & El TL431 estabiliza cuando su pin de referencia llega a $2.5\,\text{V}$. \newline Fórmula: $V_{out} = 2.5 \cdot (1 + R_{superior} / R_{inferior})$. \newline Para $12\,\text{V}$, la relación $R_{superior} / R_{inferior}$ debe ser $3.8$. Con el potenciómetro se logra el ajuste fino alrededor de los $10.26\,\text{k}\Omega$. \\ \hline
        Arranque ($R_{start}$) & $\mathbf{2 \times 100\,\text{k}\Omega}$ / $1\,\text{W}$ & Resistencias de "Kick-start" para alimentar al UC384x inicialmente. \\ \hline
        Filtro EMI ($C_{Y1,2}$) & \textbf{$2.2\,\text{nF}$ / $250\,\text{VAC}$} & Capacitores de seguridad \textbf{Clase Y1}. \\ \hline
    \end{tabular}
\end{table}

\textit{Consideración térmica:} Al entregar $120\,\text{W}$, el MOSFET primario y el diodo Schottky secundario disiparán una cantidad considerable de calor (aprox. $6\,\text{W}$ a $12\,\text{W}$ en total dependiendo de la eficiencia). Es imperativo montar ambos semiconductores sobre \textbf{disipadores de aluminio} adecuados, utilizando pasta térmica y aislantes de mica o silicona si fuese necesario.

\newpage
\section{Guía de Diagnóstico y Puntos de Prueba (Troubleshooting)}

Una parte fundamental del diseño de cualquier equipo electrónico es prever cómo será diagnosticado y reparado en caso de fallo. El término en inglés para este proceso es \textit{Troubleshooting}. Para facilitar esta tarea, se marcan en el esquema puntos de prueba (TP, del inglés \textit{Test Point}) en nodos críticos del circuito. A continuación, se describen los puntos de prueba añadidos a la Figura 2 y qué se espera medir en ellos.

\begin{table}[H]
    \centering
    \caption{Puntos de Prueba (Test Points) para diagnóstico de la SMPS.}
    \renewcommand{\arraystretch}{1.3}
    \begin{tabular}{|p{1.5cm}|p{3.5cm}|p{9.5cm}|}
        \hline
        \textbf{Punto} & \textbf{Ubicación} & \textbf{Señal Esperada y Diagnóstico} \\ \hline
        \textbf{TP1} & Vbus (+) & \textbf{DC Alto Voltaje (aprox. 310V DC).} Es la tensión principal rectificada. Si no hay tensión aquí, el problema está en la etapa de entrada (Figura 1). \\ \hline
        \textbf{TP2} & Vbus (-) / Source & \textbf{Referencia de tierra primaria.} Debe estar a 0V respecto a la tierra del primario. Se usa como referencia para medir TP3, TP4 y TP5. \\ \hline
        \textbf{TP3} & Gate del MOSFET & \textbf{Onda cuadrada (aprox. 0-15V).} Es la señal de excitación del PWM. Si no hay pulsos aquí, el controlador UC384x no está oscilando o está dañado. \\ \hline
        \textbf{TP4} & Drain del MOSFET & \textbf{Onda cuadrada de alta tensión.} Es el nodo de conmutación. Se observan pulsos que van desde casi 0V hasta Vbus + el pico de \textit{flyback}. \textbf{¡PRECAUCIÓN! Medir aquí requiere una punta de osciloscopio de alta tensión.} \\ \hline
        \textbf{TP5} & Pin Isense & \textbf{Pulsos de tensión en forma de rampa (0 a 1V).} Representa la corriente que fluye por el primario. Si esta tensión es excesiva, el PWM entrará en modo de protección. \\ \hline
        \textbf{TP6} & VCC del UC384x & \textbf{DC estable (aprox. 12-16V).} Es la alimentación del controlador. Si esta tensión es baja o inestable, la fuente no arrancará o se reiniciará constantemente (hiccup mode). \\ \hline
        \textbf{TP7} & VREF del UC384x & \textbf{DC muy estable (exactamente 5V).} Es la referencia interna del controlador. Si no está en 5V, el integrado está dañado. \\ \hline
        \textbf{TP8} & Referencia del TL431 & \textbf{DC muy estable (exactamente 2.5V).} Es el corazón del lazo de realimentación. Si este voltaje no es 2.5V, la fuente no está regulando correctamente. Si es mayor, la salida es muy alta; si es menor, la salida es muy baja. \\ \hline
        \textbf{TP9} & Salida del rectificador secundario & \textbf{Pulsos DC de alta frecuencia.} Es la energía que llega desde el transformador antes de ser filtrada por el inductor $L_{out}$. \\ \hline
        \textbf{TP10} & Salida (+) & \textbf{DC estable y regulada (12V).} Es la salida final de la fuente. \\ \hline
        \textbf{TP11} & Salida (-) & \textbf{Referencia de tierra secundaria.} Debe estar a 0V respecto a la salida. \\ \hline
    \end{tabular}
\end{table}

\end{document}